% ===== §A' The Evolution of Aesthetic Judgement =====
\section{The Evolution of Aesthetic Judgement: A Dynamics in the Reproduction Cycle}
\label{p21:sec:evolution}

\noindent
The philosophy of beauty just developed treated aesthetic judgement as though it were a single event, a reconciliation achieved at a moment, an accord struck once. This is an abstraction, and it must now be corrected, for aesthetic judgement is not a point but a \emph{process}: it evolves in time, and its evolution is governed by exactly the reproduction cycle the paper analysed in its political economy. To read the same poem today and a year ago is to make two different aesthetic judgements of it; the beauty one finds in a place, a person, an experience, is shifts rather than fixed, deepens, sometimes fades, across the repeated returns of a life. This temporal evolution of aesthetic judgement is the missing dynamical dimension of the paper's aesthetics, and supplying it joins the third movement (the aesthetic) to the second (the political-economic) at their root, showing them to be two aspects of one process. The section is brief but structurally pivotal: it is the hinge on which the paper's aesthetics and its political economy turn out to be the same theory seen twice.

% ============================================================
\subsection{Aesthetic judgement is not a point but a trajectory}
\label{p21:sub:ev-trajectory}

Consider the simplest case. One reads a poem, and finds it beautiful, or does not, or finds in it a particular beauty. One returns to it a year later, and the judgement has changed: the poem one found slight now seems profound, or the poem that dazzled now seems hollow, or, most commonly, the same poem yields a different beauty, inflected by all that has happened in the intervening time, the lines that meant little now heavy with a significance that one's own life has since supplied. The aesthetic judgement of the poem is not a fixed fact about the poem-and-reader; it is a state that evolves, a point on a trajectory that the repeated returns trace out. And what is true of the poem is true of every object of aesthetic judgement: the place revisited, the music reheard, the person re-beheld, the shared experience re-membered. Aesthetic judgement is temporal through and through, a thing that develops across its returns, and to treat it as a single timeless accord is to freeze a moving thing.

This is not a defect in aesthetic judgement to be corrected toward some stable ``true'' judgement that further acquaintance approaches. The evolution is not convergence toward a fixed correct answer, as the accumulation of evidence converges toward a fact; it is genuine development, the judgement becoming \emph{other} as the judging subject and its history become other, with no terminus at which the ``real'' beauty of the thing would finally be fixed. This follows directly from the resolution of the antinomy (\xref{p21:sub:ap-relationalsubject}): if beauty is the three-register reconciliation achieved by a relational subject rather than an objective property of the thing, then as the subject evolves, as its symbolic world deepens, its imaginary investments shift, its encounters with the Real accumulate, the reconciliation it achieves with the same object evolves too. The beauty is not in the object waiting to be correctly read; it is in the accord, and the accord is struck by a subject that is itself in motion. Aesthetic judgement evolves because the relational subject that judges is is itself continually being generated and re-generated rather than static, which is to say, reproduced.

% ============================================================
\subsection{The reproduction cycle as the engine of evolution}
\label{p21:sub:ev-reproduction}

What drives the evolution? The answer joins this section to the paper's political economy: aesthetic judgement evolves through the \emph{reproduction cycle} (\xref{p21:sec:reproduction}). Recall that the reproduction of experience is the labour of shared retelling by which a generated experience is returned to the relation's present, re-felt, and re-woven into its living fabric (\xref{p21:sub:reproduction-retelling}). Each such return is a re-generation rather than a neutral replay, in which the experience is inflected by all that has happened since and added to in the re-feeling. And the aesthetic judgement of the experience is among the things so re-generated. When two people return, in retelling, to a shared journey, they do not merely recall its beauty; they \emph{re-judge} it, and the re-judging is shaped by everything the intervening time has brought, so that the beauty they find in the remembered experience evolves with each return. The reproduction cycle is therefore the engine of the evolution of aesthetic judgement: each turn of the cycle re-generates the judgement, moves it, deepens or alters it, traces another point on its trajectory.

This is why the example with which this paper's political economy was concerned, \emph{do you remember that time we got lost}, is at once a reproduction of experience and an evolution of aesthetic judgement. The getting-lost, judged in the moment as a frustration, is re-judged in the retelling, and across repeated retellings its aesthetic value evolves: what was experienced as an irritation becomes, through the reproductive returns, perceived as a charm, then as a treasured oddity, then as one of the most beautiful moments of the whole journey. The aesthetic judgement of the experience has traversed a trajectory, frustration to charm to treasure, and the reproduction cycle is what carried it along that path. This is the precise mechanism by which the conflict analysed in the previous section (\xref{p21:sub:conflict-wealth}) can become wealth: the reproduction cycle evolves the aesthetic judgement of the conflict from the negative value it had in the moment to the positive value it acquires in the transfigured retelling. The evolution of aesthetic judgement and the reproduction of experience are not two processes but one: to reproduce an experience is to evolve its aesthetic judgement, and the trajectory of the judgement is the track the reproduction cycle traces.

% ============================================================
\subsection{The dynamics: good and bad evolution}
\label{p21:sub:ev-dynamics}

If aesthetic judgement evolves through the reproduction cycle, then it has a \emph{dynamics}, and the dynamics can run well or badly, which connects this section to the good and bad cycles the series has analysed throughout, and to the dynamical layer of its formal companion \parencite{paperix,wan2026braid}. The evolution of aesthetic judgement is the dynamical state of a system, and like any such system it can move toward enrichment or toward collapse.

The good evolution is the deepening one. Here each turn of the reproduction cycle adds to the aesthetic judgement, so that the beauty found in the experience grows richer with each return, accumulating significance, gathering the resonances of all the intervening life, becoming over years an ever-deeper well of meaning. This is the good cycle in the register of aesthetic judgement: a circulation that returns more than it took, the holonomy of the judgement accumulating positive value along the path of its re-generation (\xref{p21:sub:reproduction-cycle}). The poem returned to over a lifetime that deepens with each reading, the shared journey whose beauty grows richer with each retelling, the relationship whose history becomes more beautiful as it is re-judged across the years, these are the good evolution, the aesthetic judgement enriched by its own reproduction. The bad evolution is the degrading one, and it takes two forms corresponding to the two failures of the cycle. The first is \emph{fixation}: the judgement frozen, the reproduction become rote, the same beauty found in the same way each time with no re-generation, the living judgement hardened into a fixed verdict, the solidification of the aesthetic judgement, its circulation fallen to zero \parencite{wan2026braid}. The poem reduced to a fixed ``great poem'' one no longer actually reads, the journey hardened into a set anecdote no longer re-felt: the aesthetic judgement has stopped evolving, frozen into a dead possession. The second is \emph{depletion}: the judgement worn away by a reproduction that extracts rather than enriches, the beauty consumed until nothing is left, the over-told story whose charm is exhausted, the experience re-judged so often and so carelessly that its aesthetic value is spent. Both are the bad cycle in the aesthetic register: the judgement either frozen or consumed, its living evolution arrested or its value extracted.

This gives the paper's aesthetics its dynamical completion, and it gives the cultivation of aesthetic judgement (\xref{p21:sec:aesthetic-education}) its temporal object. To cultivate aesthetic judgement is to tend its evolution over time rather than only to refine it at a moment, to keep its reproduction generative rather than rote or extractive, to return to the beautiful in ways that deepen rather than fix or deplete the judgement of it. The mutual aesthetic education of two (\xref{p21:sub:ae-cultivation}) is, seen now, the joint tending of the evolution of their shared aesthetic judgements: the labour by which two people, returning together to the beauties they have shared, keep the evolution of their common aesthetic judgement on the deepening path rather than letting it freeze or deplete. And the self-presence of relational dynamics (\xref{p21:sec:selfpresence}) is completed in the same way: the relation's dynamics that come to presence include the evolving dynamics of the two subjects' aesthetic judgements moving, through the reproduction cycle, toward an ever more deeply shared accord. The aesthetic and the political-economic, the judgement and its reproduction, the perception of beauty and the dynamics of the relation, are here seen to be one moving system, and the beauty of a long shared life is, in the end, the good evolution of a shared aesthetic judgement, deepened across a lifetime of generative reproduction. \el{χαλεπὰ τὰ καλά}: the beautiful is difficult, and now its difficulty is seen to be also temporal, the difficulty not only of striking the accord but of tending its evolution, over a whole life, on the deepening path.

% ============================================================
\subsection{The dialectic of the general and the particular}
\label{p21:sub:ev-dialectic}

The account of the evolution of aesthetic judgement toward an ``ever more deeply shared accord'' has been stated, so far, as though the goal were the convergence of two judgements into one, as though the ideal terminus were two people who judge beauty identically. This is wrong, and the correction is not a minor qualification but the introduction of the dialectic that governs the whole matter: the dialectic of the general and the particular, which is the form the materialist dialectic takes in the domain of aesthetic judgement, and which this section must now make explicit, for without it the paper's aesthetics would collapse into exactly the homogenisation its relational ontology forbids.

The problem is this. Two relational subjects cannot be required to have identical aesthetic judgements, to do so would be to annul their distinctness, to demand the fusion that this series has refused from its foundation, the merging of two into one that destroys the resonance it pretends to perfect (\xref{p21:sub:sp-generation}). A ``we'' whose two members judged all beauty identically would not be a relation of two but a single judging subject wearing two faces; the particularity of each, the irreducible difference of their perception, is a condition of there being two to relate rather than an obstacle to the relation. And yet a shared aesthetic judgement, a \emph{general} accord, a common way of finding things beautiful, is genuinely needed, for without it there is no ``we'' that judges, no shared sensibility, no common aesthetic life, only two strangers who happen to perceive in parallel. The relation requires both: a generality (the shared accord, the common sensibility) and a particularity (the irreducible difference of each judgement, the distinctness of each perceiving subject). And these two requirements pull against each other, the generality toward convergence, the particularity toward divergence, so that the matter is a \emph{dialectical} tension to be held rather than a simple good to be maximised.

The resolution is the materialist dialectic's resolution of the general and the particular, and it is precisely the concept of the \emph{concrete universal} that this paper invoked at the outset as its method (\xref{p21:sub:intro-method}). The genuine general is not the abstract universal that is reached by stripping away all particularity, the empty common residue, the homogenised judgement that all must share by having surrendered what is their own. The genuine general is the \emph{concrete} universal: the commonality that exists \emph{in and through} the particulars, generated by their relation, real only as the shared life of distinct perceivers who remain distinct. The shared aesthetic judgement of a ``we'' is general in this concrete sense: not an identity imposed upon two by the erasure of their difference, but a commonality generated between two who remain different, a shared accord that lives through their particular judgements rather than in place of them. This is the materialist insight against the idealist abstraction: the universal is not above the particulars, subsuming them; it is among them, generated by their relation, and has no existence apart from the particulars whose commonality it is \parencite{marx1992}. Applied here: the shared aesthetic judgement is not a third judgement above the two, to which both must conform, but the living commonality that the two particular judgements generate between them while remaining two. And the tension between the general and the particular, between the pull toward shared accord and the pull toward irreducible difference, is not a defect to be resolved by the victory of one side but the very \emph{motor} of the relation's aesthetic development, the contradiction whose ongoing motion is the relation's aesthetic life. The materialist dialectic holds that contradiction is the source of development; here the contradiction of the general and the particular is the source of the development of the shared aesthetic judgement, which advances not by resolving the tension but by its perpetual working.

This forbids two opposite errors, and naming them shows why the dialectic is necessary. The first error is \emph{homogenisation}: the collapse of the particular into the general, the two judgements converging into identity, the difference erased, which is the fusion the relational ontology forbids, and which, as the next subsection will show, is also a dynamical catastrophe. The second error is \emph{atomisation}: the collapse of the general into the particular, the two judgements with no commonality, no shared accord, no ``we'' that judges, which is the dissolution of the relation into two strangers. The good aesthetic life of a ``we'' lies in neither, but in the held tension of both: a shared accord (general) genuinely generated, through and not against the irreducible difference (particular) of two perceivers who remain two. It is to a light formal rendering of this dialectic, which shows, with a precision the prose alone cannot, why both the general and the particular are \emph{required} rather than merely permitted, and required for the sake of generativity itself, that the section finally turns.

% ============================================================
\subsection{A light formal rendering: two coupled fields}
\label{p21:sub:ev-field}

The dialectic of the general and the particular can be given a formal rendering, and the rendering earns its place by showing something the prose can assert but not demonstrate: that the particularity of the two judgements is \emph{necessary for the survival of generativity itself} rather than merely tolerable. The rendering is offered in the spirit the whole paper has maintained toward the formal, lightly, as a dialectical suggestion and not a formal assertion, for the beautiful is difficult and a true account of it must preserve its openness rather than close it in a model. What follows is a gesture toward a structure, not a theory of it; its value is heuristic, the lending of a precise image to a dialectical thought, and it is to be read as such.

Picture the two subjects' aesthetic judgements as two fields, $\phi_1$ and $\phi_2$, each a configuration evolving in time through the space of possible judgements, each the trajectory of one subject's sensibility. Picture their evolution as governed by an action, or equivalently a dynamics, with two kinds of term. There is a \emph{coupling} term, write it schematically as a cost that penalises the divergence of the two fields, something of the form $\kappa\,V(\phi_1 - \phi_2)$, which pulls the two judgements toward agreement and is minimised when they accord. And there is, for each subject, a \emph{self} term, a proper dynamics $f_i(\phi_i)$ of each field on its own, the irreducible individual tendency of each subject's sensibility, which owes nothing to the other and follows its own form. The coupling term is the \emph{general}: it is what draws the two judgements toward a shared accord, and it is what makes possible their joint, gradual evolution toward a common attractor, a shared aesthetic sensibility, a basin of agreement into which the two can settle together. The self term is the \emph{particular}: it is what keeps each judgement its own, what preserves the dynamical diversity of the two fields, what prevents their collapse into a single configuration.

The dialectic is now legible in the dynamics, and it yields exactly the two errors the previous subsection named, plus the deep reason the particular is required. A system of \emph{pure coupling}, the self terms suppressed, only the divergence-penalising term left, collapses: the two fields are driven to identity, settling into a single shared configuration, all difference annihilated. This is homogenisation, and in the dynamics it is a positive danger rather than merely a loss of richness, for a system collapsed to a single homogeneous configuration has lost all dynamical diversity, and a system without diversity \emph{cannot escape a degenerating attractor}. Should the shared attractor into which the two judgements have collapsed begin to degenerate, to become a ``bad'' attractor, a fixed and frozen configuration, the solidification analysed above (\xref{p21:sub:ev-dynamics}), the homogenised system has no resource by which to leave it: both fields are at the same dead point, and there is no perturbation, no divergent tendency, no alternative trajectory to carry the system out. The generativity of the system is extinguished, frozen at a dead attractor with no path away. This is the dynamical meaning of the fusion the relational ontology forbids: not merely the loss of two-ness, but the loss of the diversity that is the condition of continued evolution, and so the extinction of generativity.

And here is the deep role of the particular, the self term, the irreducible individuality of each judgement: \emph{it is the possibility of escape from the extinction of generativity}. Because the two fields retain their own proper dynamics, their own divergent tendencies, the system retains diversity, and diversity is precisely what allows a system to leave a degenerating attractor. When the shared attractor begins to die, the self term of one or the other field can carry it off the dead point, perturb the system, open a divergence that provides a direction of escape, and so migrate the system toward a new attractor where generativity can resume. The particularity of the two judgements is, in this light, not a romantic concession to individuality but a structural necessity for the survival of the relation's generativity: it is the diversity reserve, the evolutionary variance, the source of the perturbations by which a relation escapes the death of its shared accords. A relation whose two members have become aesthetically identical has spent this reserve and is, for that reason, fragile, unable to renew a sensibility that has gone stale, trapped in whatever shared judgement it has frozen into. A relation whose two members remain aesthetically distinct, while genuinely sharing, retains the reserve, and can therefore keep escaping the deaths of its accords, keep migrating to new and living attractors, keep its generativity from extinction. The particular is required \emph{for the sake of generativity}, which is the relation's own ultimate end \parencite{paperix}; and so the dialectic of the general and the particular is a unity in which each serves the relation's generativity rather than a standoff between two goods, the general by giving the two a shared aesthetic life, the particular by keeping that life able to renew itself.

This is the materialist dialectic in its proper form: not the triumph of the general over the particular (homogenisation, the dead shared attractor) nor of the particular over the general (atomisation, no shared life at all), but their contradictory unity, whose perpetual motion is the relation's living aesthetic development, and whose two poles are each necessary, the general for the shared accord, the particular for the escape from the accord's death. The coupling builds the ``we'''s common sensibility; the self terms keep it alive; and the good evolution of a shared aesthetic judgement is the dialectical motion of the two, a coupled system that settles toward shared accords generative enough to be worth sharing and diverse enough to be escaped when they die. The model is, as promised, only a gesture, a lending of dynamical precision to a dialectical thought, not a closing of the matter into a formalism, for the matter does not close. But it shows, with a clarity the prose alone could not reach, why the paper's relational aesthetics requires both the general and the particular, and why their tension is the very life of a shared aesthetic existence rather than a problem to be solved. \el{χαλεπὰ τὰ καλά}: the beautiful is difficult, and its difficulty is, here, the difficulty of holding a contradiction in living motion, of being, with another, neither identical nor separate, but two who share a beauty that remains, in each, their own.
